% GENERATED FILE. Edit canonical lessons, then run npm run generate:books.
% canonical-source-hash: fnv1a64:d9c601b54fd3551a
% canonical-lessons: LA-C330-funus, LA-C330-sepelio, LA-C330-sepultura, LA-C330-rogus, LA-C330-exsequiae, LA-R330-review-323-330

\chapter{Fūnus}
\label{ch:la-funus}
\hlchaptermodality{330}

\begin{chapteropening}
\textbf{By the end of this chapter:} \emph{I can describe a Roman funeral: the funeral, to bury, the burial, the pyre and the funeral procession.}

Complete the last lesson of chapter 330: I can describe a Roman funeral: the funeral, to bury, the burial, the pyre and the funeral procession.
\end{chapteropening}

\section[fūnus, fūneris]{fūnus, fūneris — a funeral; death}

\label{lesson:LA-C330-funus}



\begin{quote}
Before the new one: say the Latin for to die, then the Latin for to depart.
\end{quote}

\subsection*{You'll want to know: fūnus, fūneris}
\textbf{fūnus, fūneris} — \textquotedblleft{}a funeral; death\textquotedblright{}.

\begin{quote}
Fūnus māne per forum dūcēbātur.
\end{quote}

\begin{quote}
\emph{In the morning the funeral was being led through the forum.}
\end{quote}

Neuter, 3rd declension. The funerals of great men passed through the forum with music. English \emph{funeral}.

\begin{cousinweb}
\textbf{fūnus} is a Latin word whose origin is not settled; Lewis \& Short link it, with doubt, with a word for smoke. English \emph{funeral} comes from its family.
\end{cousinweb}

\begin{grammarlens}[title={What you've built}]
One more word for telling, explaining and describing at B1.
\end{grammarlens}

\subsection*{Your turn}
\begin{itemize}
\raggedright
  \item \emph{Say it:} \emph{fūnus, fūneris}
  \item \emph{Say it:} \emph{fūnus, fūneris}, once more
  \item \emph{Say it:} a sentence of your own with \emph{fūnus, fūneris}
  \item \emph{Recall:} say the Latin for to die, then the Latin for to depart, then say \emph{fūnus, fūneris} again
\end{itemize}

\begin{tcolorbox}[breakable,colback=teal!4,colframe=teal!35!black,title={Before you move on}]
What does fūnus, fūneris mean? (\textquotedblleft{}A funeral; death\textquotedblright{}.) Say it once more.
\end{tcolorbox}
\section[sepeliō, sepelīre]{sepeliō, sepelīre — to bury}

\label{lesson:LA-C330-sepelio}



\begin{quote}
Before the new one: say the Latin for to depart, then the Latin for a funeral.
\end{quote}

\subsection*{You'll want to know: sepeliō, sepelīre}
\textbf{sepeliō, sepelīre} — \textquotedblleft{}to bury\textquotedblright{}.

\begin{quote}
Mīlitēs sociōs suōs prope viam sepelīvērunt.
\end{quote}

\begin{quote}
\emph{The soldiers buried their comrades near the road.}
\end{quote}

\emph{sepeliō, sepelīre, sepelīvī, sepultum}. Graves lay outside the city walls, along the roads.

\begin{cousinweb}
\textbf{sepeliō} is an inherited verb whose earlier origin is not settled (Lewis \& Short, EtymDB). English \emph{sepulchre} belongs to its family.
\end{cousinweb}

\begin{grammarlens}[title={What you've built}]
One more word for telling, explaining and describing at B1.
\end{grammarlens}

\subsection*{Your turn}
\begin{itemize}
\raggedright
  \item \emph{Say it:} \emph{sepeliō, sepelīre}
  \item \emph{Say it:} \emph{sepeliō, sepelīre}, once more
  \item \emph{Say it:} a sentence of your own with \emph{sepeliō, sepelīre}
  \item \emph{Recall:} say the Latin for to depart, then the Latin for a funeral, then say \emph{sepeliō, sepelīre} again
\end{itemize}

\begin{tcolorbox}[breakable,colback=teal!4,colframe=teal!35!black,title={Before you move on}]
What does sepeliō, sepelīre mean? (\textquotedblleft{}To bury\textquotedblright{}.) Say it once more.
\end{tcolorbox}
\section[sepultūra, sepultūrae]{sepultūra, sepultūrae — burial}

\label{lesson:LA-C330-sepultura}



\begin{quote}
Before the new one: say the Latin for a funeral, then the Latin for to bury.
\end{quote}

\subsection*{You'll want to know: sepultūra, sepultūrae}
\textbf{sepultūra, sepultūrae} — \textquotedblleft{}burial\textquotedblright{}.

\begin{quote}
Sepultūra in sepulcrō familiae facta est.
\end{quote}

\begin{quote}
\emph{The burial took place in the family tomb.}
\end{quote}

Feminine, from \emph{sepultum}, the supine of \emph{sepeliō}. English \emph{sepulture}.

\begin{cousinweb}
\textbf{sepultūra} is built on \emph{sepultus}, the participle of \emph{sepeliō}, \textquotedblleft{}to bury\textquotedblright{} (EtymDB), with the ending \emph{-ūra}.
\end{cousinweb}

\begin{grammarlens}[title={What you've built}]
One more word for telling, explaining and describing at B1.
\end{grammarlens}

\subsection*{Your turn}
\begin{itemize}
\raggedright
  \item \emph{Say it:} \emph{sepultūra, sepultūrae}
  \item \emph{Say it:} \emph{sepultūra, sepultūrae}, once more
  \item \emph{Say it:} a sentence of your own with \emph{sepultūra, sepultūrae}
  \item \emph{Recall:} say the Latin for a funeral, then the Latin for to bury, then say \emph{sepultūra, sepultūrae} again
\end{itemize}

\begin{tcolorbox}[breakable,colback=teal!4,colframe=teal!35!black,title={Before you move on}]
What does sepultūra, sepultūrae mean? (\textquotedblleft{}Burial\textquotedblright{}.) Say it once more.
\end{tcolorbox}
\section[rogus, rogī]{rogus, rogī — a funeral pyre}

\label{lesson:LA-C330-rogus}



\begin{quote}
Before the new one: say the Latin for to bury, then the Latin for burial.
\end{quote}

\subsection*{You'll want to know: rogus, rogī}
\textbf{rogus, rogī} — \textquotedblleft{}a funeral pyre\textquotedblright{}.

\begin{quote}
Fīlius rogum patris accendit.
\end{quote}

\begin{quote}
\emph{The son lit his father's pyre.}
\end{quote}

Masculine. The pyre was built of wood outside the walls; the ash was gathered into an \emph{urna} (taught).

\begin{cousinweb}
\textbf{rogus} is inherited; EtymDB links it with a Proto-Indo-European root, but the link is not settled.
\end{cousinweb}

\begin{grammarlens}[title={What you've built}]
One more word for telling, explaining and describing at B1.
\end{grammarlens}

\subsection*{Your turn}
\begin{itemize}
\raggedright
  \item \emph{Say it:} \emph{rogus, rogī}
  \item \emph{Say it:} \emph{rogus, rogī}, once more
  \item \emph{Say it:} a sentence of your own with \emph{rogus, rogī}
  \item \emph{Recall:} say the Latin for to bury, then the Latin for burial, then say \emph{rogus, rogī} again
\end{itemize}

\begin{tcolorbox}[breakable,colback=teal!4,colframe=teal!35!black,title={Before you move on}]
What does rogus, rogī mean? (\textquotedblleft{}A funeral pyre\textquotedblright{}.) Say it once more.
\end{tcolorbox}
\section[exsequiae, exsequiārum]{exsequiae, exsequiārum — a funeral procession, funeral rites}

\label{lesson:LA-C330-exsequiae}



\begin{quote}
Before the new one: say the Latin for burial, then the Latin for a funeral pyre.
\end{quote}

\subsection*{You'll want to know: exsequiae, exsequiārum}
\textbf{exsequiae, exsequiārum} — \textquotedblleft{}a funeral procession, funeral rites\textquotedblright{}.

\begin{quote}
Multī cīvēs exsequiās cōnsulis secūtī sunt.
\end{quote}

\begin{quote}
\emph{Many citizens followed the consul's funeral procession.}
\end{quote}

Feminine, plural: those who follow the body out. English \emph{exequies}.

\begin{cousinweb}
\textbf{exsequiae} is built on \emph{exsequor}, \textquotedblleft{}to follow out\textquotedblright{} (Lewis \& Short), from \emph{ex} and \emph{sequor}, \textquotedblleft{}to follow\textquotedblright{}: the following of the dead.
\end{cousinweb}

\begin{grammarlens}[title={What you've built}]
One more word for telling, explaining and describing at B1.
\end{grammarlens}

\subsection*{Your turn}
\begin{itemize}
\raggedright
  \item \emph{Say it:} \emph{exsequiae, exsequiārum}
  \item \emph{Say it:} \emph{exsequiae, exsequiārum}, once more
  \item \emph{Say it:} a sentence of your own with \emph{exsequiae, exsequiārum}
  \item \emph{Recall:} say the Latin for burial, then the Latin for a funeral pyre, then say \emph{exsequiae, exsequiārum} again
\end{itemize}

\begin{tcolorbox}[breakable,colback=teal!4,colframe=teal!35!black,title={Before you move on}]
What does exsequiae, exsequiārum mean? (\textquotedblleft{}A funeral procession, funeral rites\textquotedblright{}.) Say it once more.
\end{tcolorbox}
\section[(dialogue)]{Words from chapters 323–330 — a review}

\label{lesson:LA-R330-review-323-330}



\begin{quote}
Say the Latin for to give birth to and a birth.
\end{quote}

\subsection*{Your turn}
Say these aloud:

\begin{itemize}
\raggedright
  \item to give birth to, a birth, a midwife, to suckle, a nurse
  \item a cradle, a boy, infancy, childhood, to bring up
  \item an amulet, a rattle, a first name, an extra family name, a guardian
  \item young adulthood, youth, old age, orphaned, a ward
  \item a bride-to-be, a betrothal, to promise in marriage, a bride, marriage
  \item a bridal veil, a matron of honour, of a wedding, a married woman, a divorce
  \item death, dead, to breathe one's last, to die, to depart
  \item a funeral, to bury, burial, a funeral pyre, a funeral procession
\end{itemize}

\begin{tcolorbox}[breakable,colback=teal!4,colframe=teal!35!black,title={Before you move on}]
To give birth to? (\textbf{pariō, parere}.) A birth? (\textbf{partus, partūs}.) A midwife? (\textbf{obstetrīx, obstetrīcis}.) To suckle? (\textbf{nūtriō, nūtrīre}.) A nurse? (\textbf{nūtrīx, nūtrīcis}.) A cradle? (\textbf{cūnae, cūnārum}.) A boy? (\textbf{puer, puerī}.) Infancy? (\textbf{īnfantia, īnfantiae}.) Childhood? (\textbf{pueritia, pueritiae}.) To bring up? (\textbf{ēducō, ēducāre}.) An amulet? (\textbf{bulla, bullae}.) A rattle? (\textbf{crepundia, crepundiōrum}.) A first name? (\textbf{praenōmen, praenōminis}.) An extra family name? (\textbf{cognōmen, cognōminis}.) A guardian? (\textbf{tūtor, tūtōris}.) Young adulthood? (\textbf{adulēscentia, adulēscentiae}.) Youth? (\textbf{iuventūs, iuventūtis}.) Old age? (\textbf{senectūs, senectūtis}.) Orphaned? (\textbf{orbus, orba, orbum}.) A ward? (\textbf{pūpillus, pūpillī}.) A bride-to-be? (\textbf{spōnsa, spōnsae}.) A betrothal? (\textbf{spōnsālia, spōnsālium}.) To promise in marriage? (\textbf{dēspondeō, dēspondēre}.) A bride? (\textbf{nūpta, nūptae}.) Marriage? (\textbf{mātrimōnium, mātrimōniī}.) A bridal veil? (\textbf{flammeum, flammeī}.) A matron of honour? (\textbf{prōnuba, prōnubae}.) Of a wedding? (\textbf{nūptiālis, nūptiāle}.) A married woman? (\textbf{mātrōna, mātrōnae}.) A divorce? (\textbf{dīvortium, dīvortiī}.) Death? (\textbf{mors, mortis}.) Dead? (\textbf{mortuus, mortua, mortuum}.) To breathe one's last? (\textbf{exspīrō, exspīrāre}.) To die? (\textbf{obeō, obīre}.) To depart? (\textbf{dēcēdō, dēcēdere}.) A funeral? (\textbf{fūnus, fūneris}.) To bury? (\textbf{sepeliō, sepelīre}.) Burial? (\textbf{sepultūra, sepultūrae}.) A funeral pyre? (\textbf{rogus, rogī}.) A funeral procession? (\textbf{exsequiae, exsequiārum}.)
\end{tcolorbox}
